% problem-000051 3 元素发现与硫化物转化
\section*{3. 元素发现与硫化物转化}
将含元素A的硫化物矿在 $\mathrm { N } _ { 2 }$ ⽓氛中 $8 0 0 ~ ^ { \circ } \mathrm { C }$ 处理，分解产物中有A的硫化物B；随后升温⾄ $8 2 5 ~ ^ { \circ } \mathrm { C }$ 并向体系中同⼊氨⽓，得到红⾊化合物C，C溶于发烟硝酸得到⽩⾊沉淀D。经过滤洗涤，D在 $6 0 0 ~ ^ { \circ } \mathrm { C }$ 与$\mathrm { C O C l } _ { 2 }$ 反应，产物冷却后得到液体E，E遇⽔⽣成D，在E的 $6 \mathrm { m o l } \mathrm { L } ^ { - 1 }$ 盐酸溶液中通⼊ $\mathrm { H } _ { 2 } \mathrm { S }$ 得到沉淀 $\mathbf { B } ;$ ；将D溶于NaOH溶液，⽤硝酸调节H+浓度⾄约0.3 mol L−1，加⼊钼酸铵溶液常温下反应产⽣橙⻩⾊沉淀F，F与⼗⼆钼酸结构等同；将D加⼊ $\mathrm { H } _ { 3 } \mathrm { P O } _ { 2 }$ 和 $\mathrm { { \bar { { d } } _ { 3 } P O _ { 3 } } }$ 的混合溶液可得到⻩绿⾊的亚磷酸盐沉淀G，G在碱性溶液中转换为⻩⾊沉淀H，H放置时脱⽔变成I，I也可由D和A的单质在⾼温下反应产⽣，D变为I失重15.3%。

\subsection*{3-1}
写出A\~I的化学式。

\subsection*{3-2}
写出B与氨⽓反应⽣成C的反应⽅程式。

\subsection*{3-3}
写出 D 在 $\mathrm { H } _ { 3 } \mathrm { P O } _ { 2 }$ 和 $\mathrm { H _ { 3 } P O _ { 3 } }$ 中⽣成G的反应⽅程式。
